\problemname{Tvåpotensernas kraft}

Louise och hennes nära vän Elin har alltid varit helt ostoppbara tillsammans.
De två kallade sin oslagbara duo-energi för Vänskapens kraft, och när de jobbade ihop kunde ingenting stoppa dem.

Men en dag upptäckte Louise något ännu större. Något så stort att det möjligtvist inte skulle få plats i ett 32-bitars heltal.

Hon satt och räknade framför lägerelden och märkte att vissa tal växte på ett speciellt sätt: $1, 2, 4, 8, 16, 32, \dots$.  
Varje tal var dubbelt så stort som det förra, vilket gör att talen växer väldigt snabbt.

\begin{align*}
& \text{''Det här... det är mer än vänskapens kraft'', viskade Louise.} \\
& \text{''Det är tvåpotensernas kraft''.}
\end{align*}

Ja. Tvåpotensernas kraft är verkligen något speciellt. Men nu undrar Louise och Elin, ifall ett givet tal $N$ kommer någonsin dyka upp i sekvensen av tvåpotenser?

\section*{Indata}
Indatan består av ett enda heltal, $N$ ($1 \leq N \leq 10^{18}$).

\section*{Utdata}
Skriv ut \verb|Yes| om $N$ är en tvåpotens, annars skriv ut \verb|No|.


\section*{Poängsättning}
Din lösning kommer att testas på en mängd testfallsgrupper.
För att få poäng för en grupp så måste du klara alla testfall i gruppen.

\noindent
\begin{tabular}{| l | l | p{12cm} |}
  \hline
  \textbf{Grupp} & \textbf{Poäng} & \textbf{Gränser} \\ \hline
  $1$    & $20$  & $N \le 100$ \\ \hline
  $2$    & $80$  & Inga ytterligare begränsningar. \\ \hline
\end{tabular}

